Una pequeña empresa elabora cuatro tipos de zumo: manzana, piña, frutos rojos y naranja a granel.  
Cada tipo de zumo requiere cierta cantidad de fruta fresca (kg) y, salvo el de naranja a granel, botellas vacías (unidades).  
La empresa puede adquirir hasta 800 kg de fruta fresca por semana, en cualquier combinación de frutas, y dispone de exactamente 200 botellas que desea utilizar por completo para no tener sobrantes en el almacén.  

Los beneficios unitarios y consumos por litro producido son los siguientes:

\[
\begin{array}{lccc}
\text{Zumo (variable)} & \text{Fruta fresca (kg/litro)} & \text{Botellas (unidades/litro)} & \text{Beneficio (€ por litro)} \\
\hline
\text{Manzana }          & 2 & 1 & 3 \\
\text{Piña }             & 3 & 1 & 4 \\
\text{Frutos rojos }     & 4 & 1 & 5 \\
\text{Naranja a granel } & 2 & 0 & 3 \\
\hline
\end{array}
\]

La empresa desea maximizar el beneficio total, respetando el límite semanal de fruta fresca y utilizando todas las botellas disponibles.

El modelo de Programación Lineal que permite determinar el plan de producción óptimo es el siguiente:

\begin{equation}
\begin{split}
\mbox{max. } z = 3x_1 + 4x_2 + 5x_3 + 3x_4\\
s.a.:\\
2x_1 + 3x_2 + 4x_3 + 2x_4 \leq 800\\
1x_1 + 1x_2 + 1x_3 = 200\\
x_1,\,\,x_2,\,\,x_3,\,\,x_4\geq 0\\
\end{split}
\end{equation}

donde $x_1$, $x_2$, $x_3$ y $x_4$ representan, respectivamente, los litros producidos de zumo de manzana, piña, frutos rojos y naranja a granel.

Se dispone de la tabla correspondiente a una solución mediante el método de las dos fases.

\begin{center}
\begin{tabular}{c|c|cccccc|}
 & $z$ & $x_1$ & $x_2$ & $x_3$ & $x_4$ & $h_1$ & $a_2$\\ 
 & $-600$ & $0$ & $1$ & $2$ & $3$ & $0$ & $-3$ \\ 
\hline
$h_1$ & 400  & $0$ & $1$ & $2$ & $2$ & $1$ & $-2$\\ 
$x_1$ & 200  & $1$ & $1$ & $1$ & $0$ & $0$ & $1$\\ 
\hline
\end{tabular}
\end{center}



Se pide:

\begin{enumerate}[label=\alph*)]
    \item Obtener y describir el plan de producción que permita obtener beneficio máximo.

    \item Interpretar, para el plan óptimo de producción, el significado del criterio del simplex de la variable $x_2$ incluyendo en la explicación el valor de las tasas de sustitución de esa variable con respecto a las variables básicas.
    
    \item Calcular el rango dentro del cual puede variar la disponibilidad total de fruta sin que cambie la gama de zumos correspondiente a la solución óptima.

    \item Indicar cuál sería el impacto sobre el plan óptimo de producción si, en lugar de tener que utilizar obligatoriamente 200 botellas, hubiera que utilizar exactamente 199 botellas.

\end{enumerate}









 
